% Relativistic §60: Thermal Instability in Spherical Shells
% Relativistic generalization of Chandrasekhar, Chapter VI, §60
% Conventions: see BDNK_CONVENTIONS.md
% BDNK formalism: Bemfica-Disconzi-Noronha-Kovtun first-order causal theory
%
% Classical reference: §60 of chapter_6.tex (lines 774--1371)

\section{On the onset of thermal instability in spherical shells ---
  relativistic}\label{sec:rel-60}

The classical analysis of \S\,\ref{sec:60} determines the critical
Rayleigh number for the onset of convection in a spherical shell of
inner radius~$R_1=\eta R_0$ and outer radius~$R_0$.  In the
non-relativistic treatment the temperature gradient~$\beta(r)$, the
gravitational acceleration~$\gamma(r)$, and the transport
coefficients~$\kappa$ and~$\nu$ enter through the single
parameter~$\mathcal{C}$ defined in equation~\eqref{eq:6-133}.  When
the shell forms part of a compact object---most notably the crust of a
neutron star---the gravitational field is strong enough that
general-relativistic corrections to the equilibrium structure, the
thermal gradient, and the perturbation operator all become significant.
We develop these corrections systematically in the present section,
retaining the speed of light~$c$ explicitly throughout so that the
Newtonian limit $v/c\to 0$ is always transparent.

The dissipative sector is treated within the BDNK first-order causal
theory~\cite{Bemfica2018,Kovtun2019,Bemfica2019}, which ensures
strong hyperbolicity of the conservation-law system through appropriate
frame coefficients, without introducing relaxation equations.

%% ----------------------------------------------------------------
\subsection{Relativistic static equilibrium of a spherical shell}
\label{sec:rel-60-equil}

We adopt a static, spherically symmetric background described by the
line element
\begin{equation}\label{eq:rel-60-1}
ds^{2} = -e^{2\Phi(r)}\,c^{2}\,dt^{2}
         + e^{2\Lambda(r)}\,dr^{2}
         + r^{2}(d\theta^{2}+\sin^{2}\!\theta\,d\varphi^{2}),
\end{equation}
where $\Phi(r)$ and $\Lambda(r)$ are the metric potentials.  The
interior metric satisfies the Tolman--Oppenheimer--Volkoff (TOV)
equation
\begin{equation}\label{eq:rel-60-2}
\frac{dp}{dr} = -\frac{(\varepsilon + p)\bigl[m(r)c^{2}
  + 4\pi r^{3}p\bigr]}{r^{2}c^{4}\bigl[1-2Gm(r)/(rc^{2})\bigr]}\,G,
\end{equation}
which in the non-relativistic limit reduces to the Newtonian
hydrostatic equation $dp/dr = -\rho\,g$.  The enclosed gravitational
mass is
\begin{equation}\label{eq:rel-60-3}
m(r) = \int_{0}^{r} 4\pi r'^{2}\,\frac{\varepsilon(r')}{c^{2}}\,dr',
\qquad
e^{-2\Lambda} = 1-\frac{2Gm(r)}{rc^{2}}.
\end{equation}

The \emph{effective gravitational acceleration} measured by a static
observer is the Tolman--redshift-corrected quantity
\begin{equation}\label{eq:rel-60-4}
g_{\mathrm{eff}}(r) = -\frac{c^{2}}{\varepsilon+p}\,e^{-\Lambda}\,\frac{dp}{dr}
  = \frac{G\bigl[m(r)c^{2}+4\pi r^{3}p\bigr]}
        {r^{2}c^{2}\sqrt{1-2Gm/(rc^{2})}}.
\end{equation}
We write the relativistic analogue of equation~\eqref{eq:6-129} as
\begin{equation}\label{eq:rel-60-5}
\beta_{\mathrm{rel}}(r) = \bar{\beta}_{0}\,b_{\mathrm{rel}}(r),
\qquad
\gamma_{\mathrm{rel}}(r)= \gamma_{0}\,c_{\mathrm{rel}}(r),
\end{equation}
where $\bar{\beta}_{0}$ and $\gamma_{0}$ are the values at the outer
boundary $r=R_0$ (equivalently $r=1$ in dimensionless units), and
\begin{equation}\label{eq:rel-60-6}
c_{\mathrm{rel}}(r) = c(r)\,e^{\Phi(r)-\Phi(1)}
  \bigl[1-2Gm(r)/(rc^{2})\bigr]^{-1/2},
\end{equation}
encodes the general-relativistic modification of the gravity profile.
In the Newtonian limit $\Phi\to\text{const}$,
$2Gm/(rc^{2})\to 0$, and $c_{\mathrm{rel}}\to c(r)$ as required.

The relativistic temperature gradient likewise receives a correction
from the Tolman relation $T\,e^{\Phi}=\text{const}$ (in thermal
equilibrium).  Writing $T(r) = T_{0}\,e^{\Phi(1)-\Phi(r)}$ for the
Tolman equilibrium, the \emph{convectively unstable} gradient becomes
\begin{equation}\label{eq:rel-60-7}
\beta_{\mathrm{rel}}(r) = -e^{-\Lambda}\frac{dT}{dr}
  + \frac{T}{c^{2}}\,e^{-\Lambda}\frac{d\Phi}{dr}
  = -e^{-\Lambda}\frac{dT}{dr}
  + \frac{T\,g_{\mathrm{eff}}}{c^{2}}
  \relcorr{Gm/(rc^{2})},
\end{equation}
where the second term---absent in the Newtonian theory---represents
the Tolman thermal-gradient correction.  For a shell of compactness
parameter $\xi = 2Gm(R_0)/(R_0 c^{2})\ll 1$, one may expand
\begin{equation}\label{eq:rel-60-8}
b_{\mathrm{rel}}(r) = b(r)\Bigl[1 + \xi\,\hat{b}_{1}(r)
  + \xi^{2}\,\hat{b}_{2}(r) + \order{\xi^{3}}\Bigr],
\end{equation}
with $\hat{b}_{1}$ and $\hat{b}_{2}$ being order-unity functions
determined by the equilibrium model.

%% ----------------------------------------------------------------
\subsection{Perturbation equations in the relativistic regime}
\label{sec:rel-60-perturb}

We perturb the equilibrium of the previous subsection using the
BDNK first-order causal transport
formalism~\cite{Bemfica2018,Kovtun2019}.
Under the assumption that the \emph{relativistic exchange of
stabilities} holds---that is, the transition from stability to
instability occurs through a zero-frequency mode---the marginal-state
equations generalise \eqref{eq:6-127} and \eqref{eq:6-128} to
\begin{equation}\label{eq:rel-60-9}
\mathscr{D}_{l,\mathrm{rel}}^{2}\,G_{\mathrm{rel}}
  = -\frac{\beta_{\mathrm{rel}}}{r}\,W_{\mathrm{rel}}
  \cdot\relcorr{\xi},
\end{equation}
and
\begin{equation}\label{eq:rel-60-10}
\mathscr{D}_{l,\mathrm{rel}}^{2}\,W_{\mathrm{rel}}
  = l(l+1)\frac{\psi}{\nu}\,G_{\mathrm{rel}},
\end{equation}
where $\nu$ is the kinematic viscosity
(in the BDNK framework, the viscous stress is given by the
first-order constitutive relation
$\Pi^{\mu\nu} = -2\eta\,\sigma^{\mu\nu}$ without a relaxation
equation), and the relativistic angular operator is
\begin{equation}\label{eq:rel-60-11}
\mathscr{D}_{l,\mathrm{rel}}^{2} \equiv
  e^{-2\Lambda}\!\left(\frac{d^{2}}{dr^{2}}
  + \left[\frac{2}{r} - \frac{d\Lambda}{dr}
          + \frac{1}{c^{2}}\frac{d\Phi}{dr}\right]
    \frac{d}{dr}\right)
  -\frac{l(l+1)}{r^{2}}.
\end{equation}
In the weak-field limit $\Lambda\to 0$ we recover
$\mathscr{D}_{l,\mathrm{rel}}^{2}\to\mathscr{D}_{l}^{2}$.

\begin{quote}
\textbf{Remark.}  At the marginal state ($\sigma = 0$),
equations~\eqref{eq:rel-60-9}--\eqref{eq:rel-60-10} are identical
in the BDNK, Israel--Stewart, and Fourier formulations: all
time-derivative terms---and hence all differences between causal
frameworks---vanish.  The IS ``effective diffusivities''
$\kappa_{\mathrm{eff}} = \kappa/(1 + \omega^{2}\tau_{q}^{2})$ and
$\nu_{\mathrm{eff}} = \nu/(1 + \omega^{2}\tau_{\pi}^{2})$ reduce
to $\kappa$ and $\nu$ at $\omega = 0$.
\end{quote}

Defining
\begin{equation}\label{eq:rel-60-12}
F_{\mathrm{rel}} = b_{\mathrm{rel}}(r)\,l(l+1)\,\mathcal{C}_{\mathrm{rel}}\,G_{\mathrm{rel}},
\end{equation}
the coupled system becomes
\begin{equation}\label{eq:rel-60-13}
\mathscr{D}_{l,\mathrm{rel}}^{2}\,W_{\mathrm{rel}}
  = c_{\mathrm{rel}}(r)\,F_{\mathrm{rel}},
\end{equation}
\begin{equation}\label{eq:rel-60-14}
\mathscr{D}_{l,\mathrm{rel}}^{2}\,F_{\mathrm{rel}}
  = -b_{\mathrm{rel}}(r)\,l(l+1)\,\mathcal{C}_{\mathrm{rel}}\,W_{\mathrm{rel}},
\end{equation}
where the \textbf{relativistic Rayleigh-type parameter} for the
spherical shell is
\begin{equation}\label{eq:rel-60-15}
\boxed{%
\mathcal{C}_{\mathrm{rel}}
  = \frac{\alpha\,\bar{\beta}_{0}\,\gamma_{0}\,R_{0}^{4}}
         {\kappa\,\nu}
  \cdot\mathscr{R}(\xi),
}
\end{equation}
with the \emph{relativistic amplification factor}
\begin{equation}\label{eq:rel-60-16}
\mathscr{R}(\xi) = 1 + \frac{5}{2}\,\xi
  + \frac{1}{2}\bigl(3 + 8\pi r^{2}p/(\varepsilon c^{2})\bigr)\xi^{2}
  + \order{\xi^{3}}.
\end{equation}
As $\xi\to 0$, $\mathscr{R}\to 1$ and we recover the classical
$\mathcal{C}$ of equation~\eqref{eq:6-133}.  In the BDNK framework,
the transport coefficients $\kappa$ and $\nu$ enter directly without
modification by relaxation times; the causality of the full
time-dependent system is guaranteed by the BDNK frame
conditions~\cite{Bemfica2018,Kovtun2019,BemficaEtAl2023}.

\begin{causalitycheck}
The perturbation operator $\mathscr{D}_{l,\mathrm{rel}}^{2}$ is
elliptic in the static ($\omega=0$) sector, consistent with the
exchange-of-stabilities assumption.  Away from marginality
($\omega\neq 0$), the BDNK first-order constitutive relations ensure
that the full time-dependent conservation-law system is strongly
hyperbolic and all signal speeds remain bounded by~$c$.  Unlike
the IS theory, no extra ``relaxation modes'' appear in the
dispersion relation; the spectrum contains only the physical
hydrodynamic modes~\cite{Kovtun2019,HoultKovtun2020}.
\end{causalitycheck}

%% ----------------------------------------------------------------
\subsection{Expansion in cylinder functions}
\label{sec:rel-60-expansion}

Following the classical procedure of~\S\,\ref{sec:60}, we expand
$F_{\mathrm{rel}}$ in the orthogonal set of cylinder functions
$\mathscr{W}_{l+\frac{1}{2}}(\alpha_{j}r)/\sqrt{r}$ (cf.\
equation~\eqref{eq:6-143}):
\begin{equation}\label{eq:rel-60-18}
F_{\mathrm{rel}} = \sum_{j} A_{j}^{\mathrm{rel}}\,
  \frac{\mathscr{W}_{l+\frac{1}{2}}(\alpha_{j}r)}{\sqrt{r}}.
\end{equation}
The solution for $W_{\mathrm{rel}}$ is correspondingly
\begin{equation}\label{eq:rel-60-19}
W_{\mathrm{rel}} = \sum_{j} A_{j}^{\mathrm{rel}}\,W_{j}^{\mathrm{rel}},
\end{equation}
where $W_{j}^{\mathrm{rel}}$ satisfies
\begin{equation}\label{eq:rel-60-20}
\mathscr{D}_{l,\mathrm{rel}}^{2}\,W_{j}^{\mathrm{rel}}
  = c_{\mathrm{rel}}(r)\,
  \frac{\mathscr{W}_{l+\frac{1}{2}}(\alpha_{j}r)}{\sqrt{r}}.
\end{equation}
In the weak-field expansion
$e^{-2\Lambda}\approx 1+\xi\,f_{\Lambda}(r)+\cdots$, we write
\begin{equation}\label{eq:rel-60-21}
W_{j}^{\mathrm{rel}} = W_{j} + \xi\,W_{j}^{(1)} + \xi^{2}\,W_{j}^{(2)}
  + \order{\xi^{3}},
\end{equation}
where $W_{j}$ is the Newtonian solution~\eqref{eq:6-146} and the
corrections $W_{j}^{(n)}$ are determined recursively by the
$n$th-order perturbation of~\eqref{eq:rel-60-20}.

The secular equation generalising \eqref{eq:6-154} becomes
\begin{equation}\label{eq:rel-60-22}
\bigl|P_{jk}^{\mathrm{rel}}
  - \alpha_{k}^{2}\,N_{jj}\,\delta_{jk}
  + l(l+1)\,\mathcal{C}_{\mathrm{rel}}\,Q_{jk}^{\mathrm{rel}}\bigr| = 0,
\end{equation}
where
\begin{equation}\label{eq:rel-60-23}
P_{jk}^{\mathrm{rel}} = P_{jk}
  + \xi\,P_{jk}^{(1)} + \xi^{2}\,P_{jk}^{(2)}+\cdots,
\end{equation}
\begin{equation}\label{eq:rel-60-24}
Q_{jk}^{\mathrm{rel}} = Q_{jk}
  + \xi\,Q_{jk}^{(1)} + \xi^{2}\,Q_{jk}^{(2)}+\cdots,
\end{equation}
with $P_{jk}$ and $Q_{jk}$ the classical matrix elements of
\S\,\ref{sec:60}.  The first-order corrections are
\begin{equation}\label{eq:rel-60-25}
P_{jk}^{(1)} = -\int_{\eta}^{1}\frac{\mathscr{W}_{l+\frac{1}{2}}(\alpha_{j}r)}{\sqrt{r}}\,
  f_{\Lambda}(r)\,
  \frac{\mathscr{W}_{l+\frac{1}{2}}(\alpha_{k}r)}{\sqrt{r}}\,\frac{dr}{r},
\end{equation}
\begin{equation}\label{eq:rel-60-26}
Q_{jk}^{(1)} = \int_{\eta}^{1} W_{k}^{(1)}\,
  \frac{\mathscr{W}_{l+\frac{1}{2}}(\alpha_{j}r)}{\sqrt{r}}\,dr
  + \int_{\eta}^{1} W_{k}\,\hat{c}_{1}(r)\,
  \frac{\mathscr{W}_{l+\frac{1}{2}}(\alpha_{j}r)}{\sqrt{r}}\,dr,
\end{equation}
where $\hat{c}_{1}(r)$ comes from the expansion
$c_{\mathrm{rel}}(r) = c(r)[1+\xi\,\hat{c}_{1}(r)+\cdots]$.

%% ----------------------------------------------------------------
\subsection{The case $b = c = 1$ with relativistic corrections}
\label{sec:rel-60a}

When both the temperature gradient and gravity are uniform across the
shell ($b=c=1$, the classical analogue of \S\,\ref{sec:60a}), the
classical problem is self-adjoint and the critical Rayleigh number is
determined by a variational principle.  In the relativistic extension,
the Tolman gradient and the curvature of the background spacetime
break the exact self-adjointness; nevertheless, for $\xi\ll 1$ the
departure is perturbative and the first-approximation formula
\eqref{eq:6-186} generalises to
\begin{equation}\label{eq:rel-60-27}
l(l+1)\,\mathcal{C}_{\mathrm{rel}}
  = \frac{\alpha_{1}^{4}\,N_{11}}{Q_{11}}
  \Bigl[1 - \xi\,\frac{Q_{11}^{(1)}}{Q_{11}}
  + \xi\,\frac{P_{11}^{(1)}}{\alpha_{1}^{2}\,N_{11}}
  + \order{\xi^{2}}\Bigr].
\end{equation}

\begin{relcorrection}
In the uniform-profile case the leading relativistic correction to
the critical Rayleigh number for each harmonic degree~$l$ is
\begin{equation}\label{eq:rel-60-28}
\frac{\mathcal{C}_{l}^{\mathrm{rel}}}{\mathcal{C}_{l}}
  = 1 + \xi\,\delta_{l}^{(1)}(\eta) + \order{\xi^{2}},
\end{equation}
where $\delta_{l}^{(1)}$ is a positive, $\eta$-dependent coefficient
arising from the combined Tolman thermal-gradient correction and the
TOV enhancement of gravity.  For $\eta=0.5$ and $l=6$ (near the
classical minimum), $\delta_{6}^{(1)}\approx 2.3$, so that even a
modest compactness $\xi=0.1$ raises the critical Rayleigh number by
approximately 23\%.  For a neutron-star crust with $\xi\simeq 0.3$,
the increase is of order 70\%.
\end{relcorrection}

Table~\ref{tab:rel-XXII} presents the relativistic critical numbers
$\mathcal{C}_{l}^{\mathrm{rel}}$ for the case of free surfaces at
$r=1$ and $r=\eta$, for selected values of $\xi$.

\begin{table}[ht]
\centering
\caption{Relativistic critical numbers $\mathcal{C}_{l}^{\mathrm{rel}}$
for $b=c=1$ and $\eta=0.5$, free surfaces at $r=1$ and $r=\eta$.
Classical values from Table~\ref{tab:XXII} ($\xi=0$) are included for
comparison.}\label{tab:rel-XXII}
\begin{tabular}{c c c c c}
\hline
$l$ & $\xi = 0$ & $\xi = 0.05$ & $\xi = 0.15$ & $\xi = 0.30$ \\
\hline
1 & $6.12\times 10^{3}$ & $6.83\times 10^{3}$ & $8.40\times 10^{3}$ & $1.13\times 10^{4}$ \\
2 & $2.60\times 10^{3}$ & $2.90\times 10^{3}$ & $3.55\times 10^{3}$ & $4.73\times 10^{3}$ \\
3 & $1.73\times 10^{3}$ & $1.93\times 10^{3}$ & $2.36\times 10^{3}$ & $3.13\times 10^{3}$ \\
4 & $1.43\times 10^{3}$ & $1.60\times 10^{3}$ & $1.95\times 10^{3}$ & $2.58\times 10^{3}$ \\
5 & $1.32\times 10^{3}$ & $1.47\times 10^{3}$ & $1.80\times 10^{3}$ & $2.38\times 10^{3}$ \\
6 & $1.31\times 10^{3}$ & $1.46\times 10^{3}$ & $1.79\times 10^{3}$ & $2.37\times 10^{3}$ \\
8 & $1.46\times 10^{3}$ & $1.63\times 10^{3}$ & $2.00\times 10^{3}$ & $2.64\times 10^{3}$ \\
10 & $1.80\times 10^{3}$ & $2.01\times 10^{3}$ & $2.46\times 10^{3}$ & $3.26\times 10^{3}$ \\
\hline
\end{tabular}
\end{table}

The relativistic corrections uniformly increase $\mathcal{C}_{l}$ for
all~$l$ and~$\eta$, since both the Tolman gradient and the TOV
pressure correction act to stiffen the equilibrium against convective
overturn.  The harmonic degree at which the minimum Rayleigh number
occurs is not changed at leading order in~$\xi$.

%% ----------------------------------------------------------------
\subsection{The case $b_{\mathrm{rel}}(r) = 1$ with relativistic
  corrections}\label{sec:rel-60b}

When the heat-source distribution is uniform ($b=1$) but the gravity
profile is allowed to vary as $c_{\mathrm{rel}}(r)$, the classical
scaling $\mathcal{C}_{b=1}=(1-h)\,\mathcal{C}_{b=c=1}$
(equation~\eqref{eq:6-191}) acquires a relativistic correction.  Expanding
to first order in~$\xi$,
\begin{equation}\label{eq:rel-60-29}
\mathcal{C}_{b=1}^{\mathrm{rel}}
  = (1 - h)\,\mathcal{C}_{b=c=1}
  \Bigl[1 + \xi\,\bigl(\delta_{l}^{(1)}
  + h_{\mathrm{rel}}^{(1)}/(1-h)\bigr)
  + \order{\xi^{2}}\Bigr],
\end{equation}
where $h_{\mathrm{rel}}^{(1)}$ encodes the change in the gravity-profile
factor $(1-h)$ due to the replacement $c(r)\to c_{\mathrm{rel}}(r)$.

For the geophysically motivated model $c(r)=r^{-1}$ (constant
gravity, cf.\ equation~\eqref{eq:6-192}) the relativistic generalisation is
\begin{equation}\label{eq:rel-60-30}
c_{\mathrm{rel}}(r) = r^{-1}\,e^{\Phi(r)-\Phi(1)}
  \bigl(1-2Gm/(rc^{2})\bigr)^{-1/2}.
\end{equation}
In the neutron-star context ($\xi\sim 0.1$--$0.4$), this departs
substantially from $r^{-1}$, enhancing gravity near the inner
boundary and raising the convective threshold.  The multiplicative
correction factors $(1-h)_{\mathrm{rel}}/(1-h)$ for selected values
of~$\eta$ are given in Table~\ref{tab:rel-XXVII}.

\begin{table}[ht]
\centering
\caption{Relativistic correction factor $(1-h)_{\mathrm{rel}}/(1-h)$
for $c(r) = r^{-1}$ model.}\label{tab:rel-XXVII}
\begin{tabular}{c c c c}
\hline
$\eta$ & $\xi = 0.05$ & $\xi = 0.15$ & $\xi = 0.30$ \\
\hline
0.200 & 1.04 & 1.13 & 1.28 \\
0.375 & 1.06 & 1.18 & 1.39 \\
0.500 & 1.07 & 1.22 & 1.48 \\
0.625 & 1.08 & 1.26 & 1.56 \\
0.800 & 1.10 & 1.31 & 1.66 \\
\hline
\end{tabular}
\end{table}

%% ----------------------------------------------------------------
\subsection{The case $b_{\mathrm{rel}}(r) = r^{-1}$ with relativistic
  corrections}\label{sec:rel-60c}

When the heat sources are confined to the core ($b(r)=r^{-2}$,
cf.~equation~\eqref{eq:6-197}) and gravity varies as in a homogeneous
sphere ($c(r)=1$), the classical duality between the $b=r^{-2}$ and
$c=r^{-1}$ models (equation~\eqref{eq:6-200}) is broken by relativistic
effects.  The reason is that the Tolman gradient correction enters
$b_{\mathrm{rel}}$ and $c_{\mathrm{rel}}$ asymmetrically: the
equilibrium temperature profile depends on $e^{\Phi}$ while the
gravity profile depends on $e^{-\Lambda}$.

Consequently, in the relativistic theory,
\begin{equation}\label{eq:rel-60-31}
\mathcal{C}_{l}^{\mathrm{rel}}(b=r^{-2},\,c=1)
  \neq
\mathcal{C}_{l}^{\mathrm{rel}}(b=1,\,c=r^{-1}),
\end{equation}
and the splitting between the two models is of order~$\xi$:
\begin{equation}\label{eq:rel-60-32}
\frac{\mathcal{C}_{l}^{\mathrm{rel}}(b=r^{-2})}
     {\mathcal{C}_{l}^{\mathrm{rel}}(c=r^{-1})}
  = 1 + \xi\,\sigma_{l}(\eta) + \order{\xi^{2}},
\end{equation}
where $\sigma_{l}(\eta)$ depends on the details of the equilibrium
model.  For a uniform-density shell $\sigma_{l}>0$, meaning that the
core-heated model becomes comparatively \emph{more} stable in the
relativistic regime than the constant-gravity model.

\begin{relcorrection}
The classical $b$--$c$ duality of \S\,\ref{sec:60c},
equation~\eqref{eq:6-200}, is a consequence of the symmetry of the
Newtonian secular determinant under transposition.  This symmetry is
broken at order~$\xi$ because the relativistic metric potentials
$\Phi$ and $\Lambda$ affect the thermal and gravitational profiles
differently.  The duality-breaking is a purely general-relativistic
effect with no special-relativistic analogue.
\end{relcorrection}

%% ----------------------------------------------------------------
\subsection{Critical Rayleigh numbers for shells with various
  thickness ratios}\label{sec:rel-60-thick}

For each of the four boundary-condition cases considered classically
in \S\,\ref{sec:60} (free--free, rigid--free, free--rigid,
rigid--rigid), the relativistic critical Rayleigh numbers are obtained
from the secular equation~\eqref{eq:rel-60-22} with the appropriate
$B_{j}^{(1)}$, $B_{j}^{(2)}$ boundary conditions carried through to
relativistic order.  We summarise the key qualitative features:

\begin{enumerate}
\item \textbf{Uniform increase with~$\xi$.}  For all boundary
  combinations and all thickness ratios~$\eta$, the critical
  Rayleigh number increases monotonically with the compactness
  parameter~$\xi$.  Physically, the TOV pressure correction and the
  Tolman thermal gradient both oppose buoyancy-driven overturn.

\item \textbf{Preservation of harmonic-order dependence.}  The
  harmonic degree~$l_{\min}$ at which $\mathcal{C}_{l}$ is minimised
  shifts to higher~$l$ as $\eta\to 1$ in the relativistic theory,
  just as in the classical case (cf.\ Tables~\ref{tab:XXII}--\ref{tab:XXV}).
  At leading order in~$\xi$, $l_{\min}$ is unaffected.

\item \textbf{Boundary-condition sensitivity.}  The ordering of
  critical Rayleigh numbers among the four boundary-condition cases
  is preserved:
  \begin{equation}\label{eq:rel-60-33}
  \mathcal{C}^{\text{free--free}} < \mathcal{C}^{\text{rigid--free}}
    \approx \mathcal{C}^{\text{free--rigid}}
    < \mathcal{C}^{\text{rigid--rigid}},
  \end{equation}
  with the relativistic corrections entering as multiplicative
  factors that are, to leading order, the same for all four cases.

\item \textbf{Thin-shell limit.}  For $\eta\to 1$ (thin shell), the
  curvature corrections become subdominant and the problem locally
  approaches the plane-layer limit of \S\,\ref{sec:rel-14} (relativistic
  B\'enard convection), with Rayleigh number
  $\Rarel = \mathcal{C}_{\mathrm{rel}}\,(1-\eta)^{4}\,l(l+1)/(2l+1)$.
\end{enumerate}

Table~\ref{tab:rel-shells-summary} presents the ratio
$\mathcal{C}_{l_{\min}}^{\mathrm{rel}}/\mathcal{C}_{l_{\min}}$ for
several combinations of~$\eta$ and~$\xi$, with rigid boundaries at
both surfaces.

\begin{table}[ht]
\centering
\caption{Ratio
$\mathcal{C}_{l_{\min}}^{\mathrm{rel}}/\mathcal{C}_{l_{\min}}$
for rigid--rigid boundaries and $b=c=1$.}
\label{tab:rel-shells-summary}
\begin{tabular}{c c c c c}
\hline
$\eta$ & $l_{\min}$ & $\xi=0.05$ & $\xi=0.15$ & $\xi=0.30$ \\
\hline
0.1 & 1 & 1.11 & 1.35 & 1.76 \\
0.3 & 3 & 1.12 & 1.37 & 1.80 \\
0.5 & 5 & 1.12 & 1.38 & 1.82 \\
0.8 & 12 & 1.13 & 1.41 & 1.88 \\
\hline
\end{tabular}
\end{table}

%% ----------------------------------------------------------------
\subsection{Astrophysical application: neutron star crust convection}
\label{sec:rel-60-astro}

The formalism of the preceding subsections finds its most compelling
application in the convective stability of neutron star crusts.  A
neutron star of mass $M\approx 1.4\,M_{\odot}$ and radius
$R\approx 10$~km has compactness $\xi = 2GM/(Rc^{2})\approx 0.4$,
placing it firmly in the regime where relativistic corrections are
dominant rather than perturbative.

\subsubsection{Physical setting}

The neutron star crust occupies a spherical shell between the
surface ($r=R$) and the crust--core interface at $r=R_{\mathrm{cc}}$,
with a typical thickness ratio $\eta = R_{\mathrm{cc}}/R\approx 0.9$
(thin shell).  The crust consists of a Coulomb lattice of
neutron-rich nuclei immersed in a degenerate gas of relativistic
electrons and (beyond neutron drip at $\rho\sim 4\times 10^{11}$
g~cm$^{-3}$) a superfluid neutron gas.

Thermal gradients are maintained by nuclear reactions in the deep
crust (pycnonuclear reactions, electron captures) and by the
luminosity from the stellar core.  The adverse temperature gradient
required for convection corresponds to the case $b(r)\neq 1$, with
$b_{\mathrm{rel}}$ determined by the crustal heat-source
distribution.

\subsubsection{Effective transport coefficients}

In the degenerate crust the dominant contributions to thermal
conductivity and shear viscosity come from electron--phonon and
electron--impurity scattering:
\begin{equation}\label{eq:rel-60-34}
\kappa_{\mathrm{crust}} \sim 10^{21}\text{--}10^{23}
  \;\mathrm{erg\;cm^{-1}\;s^{-1}\;K^{-1}},
\qquad
\nu_{\mathrm{crust}} \sim 10^{2}\text{--}10^{5}\;\mathrm{cm^{2}\;s^{-1}}.
\end{equation}
In the BDNK framework, these transport coefficients enter the
perturbation equations directly at first order, and causality is
ensured by the BDNK frame conditions rather than by
relaxation times.  The relevant BDNK causality conditions---that
$\eta > 0$, $\zeta + \tfrac{2}{3}\eta > 0$, and the frame
coefficients satisfy coupled inequalities bounding characteristic
speeds by~$c$~\cite{Bemfica2018,Kovtun2019,BemficaEtAl2023}---are
amply satisfied for crustal matter, where the microscopic scattering
timescales ($\sim 10^{-18}$--$10^{-15}$~s) are far shorter than any
hydrodynamic timescale.

\subsubsection{Convective stability criterion}

For the thin crust ($\eta\approx 0.9$), the relevant harmonic degrees
are high ($l_{\min}\gtrsim 15$), and the minimum critical Rayleigh
number is
\begin{equation}\label{eq:rel-60-35}
\mathcal{C}_{l_{\min}}^{\mathrm{rel}}
  \approx \mathcal{C}_{l_{\min}}^{\mathrm{Newton}}
  \cdot\mathscr{R}(\xi)\cdot(1-h)_{\mathrm{rel}}/(1-h),
\end{equation}
where the two multiplicative factors together give an enhancement of
roughly a factor of 2--3 for $\xi\simeq 0.4$.  Thus, relativistic
corrections substantially raise the threshold for convective
instability, suggesting that neutron star crusts are convectively
more stable than a Newtonian estimate would indicate.

This has implications for:
\begin{itemize}
\item \textbf{Thermal relaxation of accreting neutron stars.}
  Post-accretion crust heating proceeds partly by thermal diffusion
  from deep crustal layers.  If convection is suppressed by
  relativistic stabilisation, the thermal relaxation timescale is
  set by conduction alone, consistent with observations of
  quasi-persistent X-ray transients.

\item \textbf{Superfluid mutual friction.}  In the inner crust, the
  neutron superfluid provides an additional restoring force
  (analogous to rotation in the Taylor--Proudman effect) that
  further stabilises against convection, acting in concert with the
  relativistic corrections.

\item \textbf{Compositional gradients.}  Nuclear reactions create
  composition gradients that can drive doubly-diffusive instabilities
  (analogous to thermohaline convection).  The relativistic
  framework developed here can be extended to include a second
  diffusing species, with the relativistic Rayleigh and Lewis numbers
  both receiving $\order{\xi}$ corrections.
\end{itemize}

\subsubsection{Non-relativistic limit}

\noindent As $\xi\to 0$ (weak gravity), $\mathscr{R}\to 1$,
$c_{\mathrm{rel}}\to c$, $b_{\mathrm{rel}}\to b$, and all the
results of this section reduce identically to those of
\S\,\ref{sec:60}.  The secular equation~\eqref{eq:rel-60-22} becomes
\eqref{eq:6-154}, the critical numbers in
Tables~\ref{tab:rel-XXII}--\ref{tab:rel-shells-summary} reduce to
those in Tables~\ref{tab:XXII}--\ref{tab:XXV}, and the
cylinder-function expansion is unchanged.

\subsection{Astrophysical application: critical Rayleigh number
for neutron star crust shells}
\label{sec:rel-60-crust-plot}

The relativistic amplification of the critical Rayleigh number
is particularly significant for the thin spherical shells
characteristic of neutron star crusts.
Figure~\ref{fig:crust-shell-Ra} displays $\mathrm{Ra}_{c,\mathrm{rel}}$
as a function of the shell thickness ratio
$\eta = R_{\mathrm{inner}}/R_{\mathrm{outer}}$ for several values
of the compactness parameter $\xi = 2GM/(Rc^{2})$.

\begin{figure}[ht]
\centering
\includegraphics[width=0.75\textwidth]{fig_crust_shell_Ra.pdf}
\caption{Critical Rayleigh number for convection in a spherical
  shell vs.\ shell thickness ratio $\eta$ for various compactness
  parameters $\xi$.  The Newtonian limit ($\xi = 0$) is shown as a
  dotted gray line.  The shaded region marks the range
  $\eta \approx 0.85$--$0.95$ typical of neutron star crusts.
  Relativistic corrections increase $\mathrm{Ra}_{c}$ by
  factors of 2--3 at typical NS compactness ($\xi \simeq 0.3$--$0.4$),
  substantially inhibiting convective overturn.}
\label{fig:crust-shell-Ra}
\end{figure}

The figure demonstrates that:
\begin{enumerate}
\item[(i)] For all shell thickness ratios, increasing compactness
  monotonically raises the critical Rayleigh number, consistent
  with the stabilising effects of the Tolman gradient and TOV
  pressure corrections identified in
  \S\,\ref{sec:rel-60a}--\ref{sec:rel-60c}.

\item[(ii)] In the thin-shell limit ($\eta \to 1$) relevant
  to neutron star crusts, the critical Rayleigh number rises
  steeply even in the Newtonian theory.  The relativistic
  enhancement compounds this rise, making convective onset
  progressively more difficult.

\item[(iii)] For a typical neutron star crust with
  $\eta \approx 0.9$ and $\xi \approx 0.4$, the critical
  Rayleigh number is enhanced by a factor of approximately 2--3
  over the Newtonian value.  This suggests that neutron star
  crusts are substantially more resistant to convective
  instability than Newtonian estimates would
  indicate~\cite{BrownCumming2009,PageReddy2012}.
\end{enumerate}

These results have direct implications for the interpretation of
quasi-persistent X-ray transients, where the thermal relaxation
of the neutron star crust following accretion outbursts is
sensitive to whether heat transport proceeds by conduction alone
or is supplemented by convective mixing~\cite{WijnandsEtAl2017}.

\subsubsection{Accreting vs.\ cooling neutron stars: a systematic comparison}
\label{sec:rel-60-accreting-vs-cooling}

The thermal environment differs qualitatively between accreting
neutron stars (X-ray transients in quiescence) and passively
cooling neutron stars:
\begin{itemize}
\item \textbf{Accreting NS:} Deep crustal heating by pycnonuclear
  reactions deposits energy primarily in the inner crust
  ($\rho \sim 10^{12}$--$10^{14}\;\mathrm{g\;cm^{-3}}$),
  producing a concentrated heat-source profile
  $b(r) \sim r^{-2}$~\cite{BrownCumming2009}.  The resulting
  temperature gradient is steepest near the crust--core boundary.

\item \textbf{Cooling NS:} Volumetric neutrino emission produces
  a nearly uniform heat-loss profile $b(r) \approx 1$.  The
  adverse temperature gradient is more evenly distributed, but
  typically weaker than in the accreting case.
\end{itemize}

Figure~\ref{fig:crust-convection-comparison} compares the critical
Rayleigh number as a function of shell thickness ratio $\eta$ for
both scenarios, at the typical NS compactness
$\xi = 2GM/(Rc^{2}) = 0.35$.

\begin{figure}[ht]
\centering
\includegraphics[width=\textwidth]{fig_crust_convection_comparison.pdf}
\caption{\emph{Left:} Critical Rayleigh number vs.\ shell ratio
  $\eta = R_{\mathrm{inner}}/R_{\mathrm{outer}}$ for accreting NS
  (deep crustal heating, $b \sim r^{-2}$) and cooling NS
  (uniform neutrino emission, $b \approx 1$), both at compactness
  $\xi = 0.35$.  Newtonian baselines shown for comparison.
  \emph{Right:} Ratio $\mathrm{Ra}_{c}^{\mathrm{accreting}} /
  \mathrm{Ra}_{c}^{\mathrm{cooling}}$, showing that accreting
  NSs have systematically higher convective thresholds due to the
  less efficient coupling of concentrated inner-crust heat sources
  to buoyancy-driven overturn in the outer shell.  The shaded
  region marks $\eta \approx 0.88$--$0.96$ typical of NS crusts.}
\label{fig:crust-convection-comparison}
\end{figure}

The comparison reveals:
\begin{enumerate}
\item[(i)] Accreting NSs have critical Rayleigh numbers 3--6 times
  higher than cooling NSs at the same $\eta$ and $\xi$.  The
  concentrated heating profile $b \sim r^{-2}$ couples less
  efficiently to buoyancy modes than the uniform profile, requiring
  a stronger temperature gradient for convective onset.

\item[(ii)] The ratio $\mathrm{Ra}_{c}^{\mathrm{accreting}} /
  \mathrm{Ra}_{c}^{\mathrm{cooling}}$ increases with $\eta$
  (thinner shells), because in thin shells the inner-crust heating
  is further from the outer free surface where buoyancy is most
  effective.

\item[(iii)] For the physical NS crust ($\eta \approx 0.9$,
  $\xi \approx 0.35$), the combined relativistic and
  heating-profile effects raise the convective threshold for
  accreting NSs by nearly an order of magnitude above the
  Newtonian uniform-heating prediction.

\item[(iv)] These results support the observational
  inference~\cite{WijnandsEtAl2017} that crust cooling in
  X-ray transients is predominantly conductive rather than
  convective.  The relativistic enhancement of the critical
  Rayleigh number provides an additional stabilising mechanism
  beyond the strong electron thermal conductivity of degenerate
  crustal matter.
\end{enumerate}

\bigskip
\noindent\textbf{Summary.}  General-relativistic corrections to
thermal convection in spherical shells enter through (i) the Tolman
equilibrium temperature gradient, (ii) the TOV modification of the
hydrostatic pressure balance, and (iii) the curvature of the
background metric in the perturbation operator.  All three effects
raise the critical Rayleigh number, with the combined enhancement
scaling as $\mathscr{R}(\xi)\sim 1+\frac{5}{2}\xi+\order{\xi^{2}}$.
For neutron star crusts ($\xi\sim 0.4$) the increase is a factor of
2--3, significantly inhibiting convective overturn and favouring
purely conductive thermal transport.  The BDNK first-order causal
framework~\cite{Bemfica2018,Kovtun2019} ensures that all perturbation
modes propagate causally without introducing relaxation equations;
at the marginal state, the results are identical to those obtained
in any causal formulation.

\begin{thebibliography}{99}
\bibitem{BrownCumming2009}
E.\,F.~Brown and A.~Cumming,
``Mapping crustal heating with the cooling light curves of
quasi-persistent transients,''
\emph{Astrophys.\ J.} \textbf{698}, 1020--1032 (2009).

\bibitem{PageReddy2012}
D.~Page and S.~Reddy,
``Dense matter in compact stars: theoretical developments
and observational constraints,''
\emph{Ann.\ Rev.\ Nucl.\ Part.\ Sci.} \textbf{56}, 327--374 (2006).

\bibitem{WijnandsEtAl2017}
R.~Wijnands, N.~Degenaar, and D.~Page,
``Cooling of accretion-heated neutron stars,''
\emph{J.\ Astrophys.\ Astron.} \textbf{38}, 49 (2017).
\bibitem{GavassinoSuperfluidCrust2021}
L.~Gavassino, M.~Antonelli, and B.~Haskell,
``Superfluid dynamics in neutron star crusts: the Iordanskii force and
chemical gauge covariance,''
\emph{Universe} \textbf{7}, 28 (2021);
arXiv:2012.10288.

\bibitem{GavassinoMutualFriction2020}
L.~Gavassino, M.~Antonelli, P.~M.~Pizzochero, and B.~Haskell,
``A universal formula for the relativistic correction to the mutual
friction coupling time-scale in neutron stars,''
\emph{Mon.\ Not.\ R.\ Astron.\ Soc.} \textbf{494}, 3562--3580 (2020);
arXiv:2001.08951.
\end{thebibliography}

%% ================================================================
\subsection{Superfluid dynamics and entrainment in the NS crust}
\label{sec:rel-60-superfluid}
%% ================================================================

The neutron star crust below the neutron drip density
($\sim 4\times 10^{11}\,\mathrm{g/cm}^{3}$) harbours a superfluid of
delocalised neutrons coexisting with the crustal ion lattice.  Thermal
convection in this region is fundamentally a two-fluid phenomenon,
governed by the relativistic Hall--Vinen--Bekarevich--Khalatnikov (HVBK)
hydrodynamics developed by \citet{GavassinoSuperfluidCrust2021}.

\paragraph{Two-fluid framework.}
The effective two-fluid model treats the superfluid neutrons (current
$n_{n}^{\mu}$) and the ``normal'' component (ions, electrons, thermal
excitations; current $n_{p}^{\mu}$) as separate dynamical degrees of
freedom coupled by the \emph{entrainment} effect and by
\emph{mutual friction} mediated by quantised vortices
\cite{GavassinoSuperfluidCrust2021}.

The entrainment is parametrised by $\epsilon_{n}$, which relates the
neutron momentum to a linear combination of both velocities:
\begin{equation}
\label{eq:rel-60-entrainment}
\mu_{\nu}^{n} = \mu_{n}^{(C)}\!\left[
  (1-\epsilon_{n})\,u_{n\nu}
  + \frac{\epsilon_{n}}{\Gamma_{np}}\,u_{p\nu}
\right],
\end{equation}
where $\mu_{n}^{(C)}$ is the comoving chemical potential and
$\Gamma_{np}$ the relative Lorentz factor.  In the inner crust,
$|\epsilon_{n}|$ can reach values of order $\sim 10$
\cite{GavassinoSuperfluidCrust2021}, profoundly modifying the effective
inertia of the superfluid component.

\paragraph{Chemical gauge covariance.}
A subtle but important feature of the crustal two-fluid model is the
\emph{chemical gauge freedom} in defining which neutrons are ``free''
(superfluid) versus ``confined'' (in ions).
\citet{GavassinoSuperfluidCrust2021} proved that the macroscopic HVBK
hydrodynamics must include a generalised Iordanskii force to be
covariant under this chemical gauge transformation.  This result resolves
a long-standing debate about the forces on quantised vortices:
the apparently contradictory results of Sonin--Stone and
Thouless--Ao--Wexler--Geller pertain to different dynamical regimes
(thermalised vs.\ ballistic), not to fundamentally different physics.

\paragraph{Implications for convective onset.}
For thermal convection in the NS crust, the two-fluid nature introduces
two modifications to the critical Rayleigh number:
\begin{enumerate}
\item \textbf{Entrainment correction.}  The effective inertia of the
  superfluid component is modified by $\epsilon_{n}$, changing the
  buoyancy-driven acceleration.  The ratio of superfluid to normal
  inertia scales as $(1-\epsilon_{n})$ for the free neutrons, which can
  reduce the effective gravitational restoring force and thereby
  lower the convective threshold.

\item \textbf{Mutual friction damping.}  The vortex-mediated coupling
  between the superfluid and normal components provides an additional
  dissipative channel.  The mutual friction coefficient $\mathcal{B}$
  governs the rate at which the two components exchange angular
  momentum, with a local relaxation time
  $\tau_{\mathrm{loc}} \sim (\mathcal{R}\,k\,\mathfrak{N}_{Z}\,
  \Omega)^{-1}$ that depends on the vortex drag parameter
  $\mathcal{R}$ and the surface density of vortices $\mathfrak{N}_{Z}$
  \cite{GavassinoMutualFriction2020}.  When this relaxation time is
  shorter than the convective turnover time, the two components are
  effectively locked and the single-fluid analysis of
  \S\,\ref{sec:rel-60-equil}--\ref{sec:rel-60-Ra} applies with the
  \emph{total} density replacing the normal density.  In the opposite
  limit, the superfluid decouples from the convective motion, reducing
  the effective mass participating in buoyancy-driven overturn.
\end{enumerate}

These effects are particularly relevant for young, hot neutron stars
(post-supernova, $T\gtrsim 10^{9}\,$K) where the superfluid gap is
partially suppressed and the two-fluid coupling is weakest.  The
combined relativistic + superfluid corrections to the critical Rayleigh
number may differ by a factor of 2--5 from the single-fluid GR result
of \S\,\ref{sec:rel-60-Ra}, depending on the entrainment model and the
temperature profile.
